\section{Introduction}\label{sec:introduction}

Viterbo's conjecture states that among convex bodies of fixed volume, the
ball has the largest symplectic capacity \cite{Viterbo2000}. We work in
$\mathbb R^4$ with the Ekeland--Hofer--Zehnder capacity $c_{\rm EHZ}$: for
a convex body $K$, $c_{\rm EHZ}(K)$ is the least action of a closed
characteristic on $\partial K$, and on polytopes the minimum runs over
generalized characteristics, which follow the facet normals. In dimension
four the conjecture reads
\[
 \operatorname{sys}(K):=\frac{c_{\rm EHZ}(K)^2}{2\operatorname{vol}_4(K)}\leq1,
\]
with equality for the ball. We call $\operatorname{sys}(K)$ the systolic
ratio. It is unchanged by translations, rescaling and linear symplectic
maps.

Write $\mathbb R^4=\mathbb R^2_q\oplus\mathbb R^2_p$ with coordinates
$(q_1,q_2,p_1,p_2)$. For planar convex bodies $K_q,K_p$, the Lagrangian
product $K_q\times_L K_p$ is the set of points with $(q_1,q_2)\in K_q$ and
$(p_1,p_2)\in K_p$. The relative position of the two factors affects
capacity even though it does not affect volume. Haim--Kislev and Ostrover
disproved the conjecture with a product of two regular pentagons
\cite{HaimKislevOstrover2024}:
\[
 K_{\rm HKO}=P_5\times_L R(-\pi/2)P_5,\qquad
 \operatorname{sys}(K_{\rm HKO})=\frac{3+\sqrt5}{5}>1.
\]
Here $P_5$ is the regular pentagon of circumradius one centred at the origin
with a vertex at $(1,0)$, and $R(\theta)$ is planar rotation by $\theta$.

This thesis picks a toolbox and asks how much progress it makes on questions
around Viterbo's conjecture after this counterexample. The toolbox consists
of high-performance computation on the LICCA cluster, proof by exact
computation, standard data-science methods and AI agents. The agents are
part of the method: they ran the data science and did large parts of the
theory work. Pure mathematics has used data before: in knot theory,
patterns spotted in computed invariants have suggested results that were
then proved. Here AI both throws data science at the problem and develops
much of the theory. The questions came out of the work, and the results
answer different ones.

One of the earliest questions was whether standard data science finds new
interesting bodies with systolic ratio above one. It did not. Bodies with
systolic ratio above one appeared as small perturbations of $K_{\rm HKO}$,
none with a larger ratio; random samples and local searches from generic
starting bodies all stayed below one. The data did show a pattern: in
sampled populations, symplectic ridge area, the total symplectic area of the
two-dimensional faces normalized by volume, is strongly inversely associated
with the systolic ratio, although the association weakens among bodies
selected for low ridge area or high systolic ratio.

Because random perturbations did not improve $K_{\rm HKO}$, we conjectured
that it is a local maximum, and we prove this among polytopes with ten
facets: every sufficiently close ten-facet polytope has systolic ratio at
most $\operatorname{sys}(K_{\rm HKO})$, with equality only for translations,
rescalings and linear symplectic images. The same ten-facet space contains a
second, non-strict local maximum at systolic ratio one, a Lagrangian product
of a hexagon and a quadrilateral. Within Lagrangian products of two
independently linearly deformed regular pentagons, we prove an explicit
capacity formula and the sharp bound $(3+\sqrt5)/5$.
These proofs rest on Haim--Kislev's finite formula for the capacity of a
polytope \cite{HK2017}, its first-order variation in the facet data, and
exact arithmetic.

Section~1.1 places the EHZ capacity and the counterexample among related
results on capacities, dynamics and convex geometry. Sections~1.2--1.4
describe the local theorem, the finite capacity computations and the
data-science investigation in more detail, and Section~1.5 explains how to
read the thesis.
